\section*{Bab \ref{ch:g7:negatives} --- Bilangan Negatif}

\begin{solution}{exo:g7:negatives:1}
\omterm{def:g7:negatives:def}{Lawannya}: $-7$; \ $+3.5$; \ $0$ (\omterm{def:g7:negatives:def}{lawan} dirinya sendiri); \ $-12$. Dua
bilangan saling berlawanan jika \omterm{def:g1:addition:def}{jumlahnya} $0$.
\end{solution}

\begin{solution}{exo:g7:negatives:2}
$-4 < 1$; \quad $-4 > -7$; \quad $-2.5 < -2.4$; \quad $0 > -0.1$.
\end{solution}

\begin{solution}{exo:g7:negatives:3}
$-3.1 < -3 < -0.5 < 0 < 1 < 2.5$.
\end{solution}

\begin{solution}{exo:g7:negatives:4}
$(-5) + (-9) = -14$; \quad
$(-12) + (+7) = -5$; \quad
$(+3.5) + (-1.5) = +2$; \quad
$(-6) + (+6) = 0$.
\end{solution}

\begin{solution}{exo:g7:negatives:5}
$4 - 11 = 4 + (-11) = -7$.

$(-5) - (+3) = (-5) + (-3) = -8$.

$(-7) - (-10) = (-7) + (+10) = +3$.

$0 - (-8) = 0 + (+8) = +8$.
\end{solution}

\begin{solution}{exo:g7:negatives:6}
$(+7) + (-12) + (+3) + (-8) = (+10) + (-20) = -10$ (yang positif
bersama-sama, yang negatif bersama-sama).

$(-5.5) + (+9) + (-4.5) + (+1) = (-10) + (+10) = 0$.
\end{solution}

\begin{solution}{exo:g7:negatives:7}
$(-3) + 8 - 12 = 5 - 12 = -7$: suhunya $-7$ derajat pada malam hari.
\end{solution}

\begin{solution}{exo:g7:negatives:8}
Peletakan bebas; $E(0, -3)$ terletak pada sumbu tegak (koordinat
pertamanya $0$).
\end{solution}

\begin{solution}{exo:g7:negatives:9}
Moskwa: $(-4) - (-13) = -4 + 13 = 9$ derajat \omterm{ex:g1:subtraction:difference}{selisihnya}.
Helsinki: $2 - (-9) = 11$ derajat.
\end{solution}

\begin{solution}{exo:g7:negatives:10}
$(-12) + 7 - 4 = -5 - 4 = -9$: ia berada di $-9$ m. Untuk mencapai
permukaan (ketinggian $0$) ia harus naik $0 - (-9) = 9$ m.
\end{solution}

\begin{solution}{exo:g7:negatives:11}
$D(-2, -2)$: koordinat pertamanya sama dengan $A$, koordinat keduanya
sama dengan $C$. Sisinya: $AB = 4$ (dari $-2$ ke $2$) dan $BC = 3$
(dari $1$ turun ke $-2$), jadi \omterm{def:g6:measure:perimeter}{kelilingnya} $2 \times (4 + 3) = 14$.
\end{solution}

\begin{solution}{exo:g7:negatives:12}
Diagonal yang lengkap $0, -2, -4$ memberi \omterm{def:g1:addition:def}{jumlah} ajaib
$0 + (-2) + (-4) = -6$. Lalu, sel demi sel:
kiri-atas $= -6 - (-7) - 0 = 1$; tengah-kanan $= -6 - (-3) - (-2) =
-1$; kanan-bawah $= -6 - 0 - (-1) = -5$; tengah-bawah
$= -6 - (-7) - (-2) = 3$.
\begin{center}
\renewcommand{\arraystretch}{1.4}
\begin{tabular}{|c|c|c|}
\hline
$1$ & $-7$ & $0$ \\
\hline
$-3$ & $-2$ & $-1$ \\
\hline
$-4$ & $3$ & $-5$ \\
\hline
\end{tabular}
\end{center}
Setiap baris, kolom dan diagonalnya berjumlah $-6$.
\end{solution}

\begin{solution}{pb:g7:negatives:1}
\textbf{1.} $4\,808 - (-430) = 4\,808 + 430 = 5\,238$ m
(\cref{thm:g7:negatives:sub}: mengurangkan bilangan negatif berarti
menjumlahkan).

\textbf{2.} Saldonya:
\[
(-45) + 120 + (-60) + 15 = 75 + (-60) + 15 = 15 + 15 = 30
\text{ euro}.
\]
Untuk mengembalikan $-45$ yang semula ke $0$, setorlah \omterm{def:g7:negatives:def}{lawannya}:
$45$ euro.

\textbf{3.} $(-44) - (-100) = (-44) + 100 = 56$: Caesar hidup
$56$ tahun. (Tidak ada perangkap di sini: kedua tahunnya berada di
sisi batas era yang sama.)

\textbf{4.} Tahun $0$: membilang dari $-63$ naik ke $+14$ melewati
$0$, tetapi kalender para sejarawan melompat langsung dari
$1$ SM ke $1$ M --- tahun $0$ tidak pernah ada, jadi hitungan yang
naif itu satu tahun terlalu panjang.

\textbf{5.} $63$ SM menjadi $-(63 - 1) = -62$. Lalu
$14 - (-62) = 14 + 62 = 76$ tahun: tepat angka milik para sejarawan.
Pada garis yang dinamai ulang oleh para astronom tidak ada celah,
jadi pengurangan biasa memberi lama waktu yang benar.

\textbf{6.} $1$ SM $\to 0$; \ $10$ SM $\to -9$; \ tahun M
$476 \to +476$; \ $753$ SM $\to -752$.

\textbf{7.} $476 - (-752) = 476 + 752 = 1\,228$ tahun.

\textbf{8.} $12$ SM adalah $-11$. Kemunculan berikutnya:
$-11 + 76 = 65$, lalu $141$, lalu $217$: yaitu tahun $65$ M, $141$
M dan $217$ M. Sebelumnya: $-11 - 76 = -87$, yaitu tahun
$87 + 1 = 88$ SM milik sejarawan.

\textbf{9.} $5$ SM adalah $-4$, jadi $5 - (-4) = 9$ tahun --- bukan
$10$. Tahun \omterm{def:g1:counting:zero}{nol} yang hilang itu menyerang setiap kali lama waktunya
melintasi batas era.

\textbf{10.} $3$ SM adalah $-2$; sepuluh tahun kemudian adalah
$-2 + 10 = +8$: anak itu berumur sepuluh tahun pada tahun $8$ M.

\textbf{11.} Antara $-7$ dan $-2$: $(-2) - (-7) = 5$. Antara
$-3.5$ dan $4.5$: $4.5 - (-3.5) = 8$.

\textbf{12.} Mengelompokkan \omterm{def:g7:negatives:def}{lawan} dan kawan-kawannya:
\[
3 + (-5) + 2 + (-1) + 4
= \bigl(3 + 2\bigr) + (-5) + (-1) + 4
= 0 + (-1) + 4 = +3 :
\]
liftnya berakhir di lantai $3$. Urutannya tidak berpengaruh: \omterm{def:g1:addition:def}{jumlah}
\omterm{def:g7:negatives:def}{bilangan bertanda} dapat disusun ulang dengan bebas (setiap gerakan
ditempuh sepanjang garis yang sama, dan seluruh jalannya sama).

\textbf{13.} Perubahan bersihnya:
$2 + (-4) + 1 + 0 + (-3) + 5 + (-2) = -1$ derajat. Mulai dari
$-2$: suhu akhirnya $-2 + (-1) = -3$ derajat. Secara umum,
nilai akhirnya adalah nilai awal ditambah \omterm{def:g1:addition:def}{jumlah} semua
perubahannya --- tidak perlu mengikuti naik-turunnya satu per satu.

\textbf{14.} Kembali ke titik awal berarti perubahan bersihnya $0$:
naik dan turunnya harus meniadakan dengan tepat. Kelima perubahan
yang diketahui berjumlah $4 + (-7) + 2 + 3 + (-5) = -3$, jadi
perubahan keenamnya adalah \omterm{def:g7:negatives:def}{lawannya}: $+3$ (tiga ratus meter naik).

\textbf{15.} Menurut para astronom: $509$ SM $\to -508$, dan
$476 - (-508) = 984$ tahun. Hitungan sang sejarawan
$509 + 476 = 985$ membilang satu tahun terlalu banyak --- sekali lagi
tahun \omterm{def:g1:counting:zero}{nol} yang hantu itu, yang diam-diam masuk ke penjumlahannya
tetapi tidak ada dalam sejarah. (Sebagai pemeriksaan, kedua eranya:
Republik $-508$ sampai $-26$, yaitu $482$ tahun; Kekaisaran $-26$
sampai $+476$, yaitu $502$ tahun; $482 + 502 = 984$.)

\end{solution}
